% Canonical Paper I source.
This paper has established the syntactic and projective foundation of the
programme.  The organizing principle is that an expression should not be
identified prematurely with its endpoint value.  An unmarked sequential tree
requires a choice of accumulator; the resulting fully marked planar tree is
equivalent to a bounded marked history.  From that history, operator evaluation
and the charge shadow form separate branches, while an initial value turns an
operator into an endpoint.  The non-bijective forgetful maps of
\eqref{eq:process-result-hierarchy} lose information, and the structures proved
here are the ones that survive at the syntactic and operator levels.

\subsection{Results proved here}

The first part of the paper classified binary expression trees with a unique
legal internal evaluation order and represented them by marked spinal
histories.  Operand-slot data made it possible to distinguish planar mirror,
temporal reversal, and path inverse.  These operations act at different levels
and need not induce the same operator comparison.

For nondegenerate bilateral one-hole contexts, projective continuation produced
M\"obius transformations.  The elementary contexts generate
$\operatorname{PGL}_2(K)$ over an arbitrary field, while ordinary addition and
nonzero scaling form the affine, or Borel, sector stabilizing the chosen point
at infinity.  This placement is important: projective semantics extends the
operator domain, but it does not turn a pole into an admissible ordinary
arithmetic operation.

The smaller alphabet consisting of translation by $\pm\sqrt2$ and negative
inversion realizes the $q=4$ Hecke group inside this projective semantics.  Its
matrix relation $(JT)^4=1$ gives an explicit pair of different marked histories
with the same operator.  The subsequent factorization through cell stabilizers
separates four levels that are needed downstream: history, operator, cell coset,
and endpoint.  In particular, arithmetic words are not identified one-to-one
with tiles in a chosen Hecke cellulation.

The second part of the paper developed the zero theory of regular arithmetic
expression spaces, using only the imported interface of
Section~\ref{sec:regular-aes-interface}.  At a zero of a regular assignment one
has $|\nabla a|=|\mu|\ne0$, so the interior zero locus is a smooth
codimension-one set and cannot have isolated points, crossings, or branch
points.  Under the additional assumptions that the metric is complete and the
manifold is connected and boundaryless, unit-gradient rectification gives the
diffeomorphism
\[
  M\cong Z(a)\times\mathbb R ,
\]
hence the zero locus is nonempty and connected.  This is not asserted for
incomplete models or manifolds with boundary, and it is not in general a
Riemannian product decomposition.  The splitting theorem is therefore a no-go
statement for branching zero networks in the regular category, and the paper
accordingly defines singular arithmetic expression spaces by a closed
nowhere-dense singular set admitting no local regular extension.  One
isolated-zero model is verified in detail; a smooth parameter family with
spatially regular zeros has a submersive total zero set, and the properness
warning records exactly how far that statement may be used.

The proved logical spine may therefore be summarized as
\[
\begin{gathered}
    \text{sequential tree}
      \longrightarrow \text{marked planar sequential tree}
      \longleftrightarrow \text{bounded marked history},\\
    \text{marked history}
      \longrightarrow \text{projective operator}
      \longrightarrow \text{ordinary endpoint or pole},\\
    \operatorname{PGL}_2
      \supset \operatorname{Aff}(1)
      \longrightarrow \text{Bruhat placement},\\
    \operatorname{Hist}_4
      \longrightarrow G_4
      \longrightarrow G_4/H_{\mathcal C}
      \longrightarrow \text{Hecke cell},\\
    \text{regular AES}
      \longrightarrow Z(a)\text{ smooth}
      \longrightarrow M\cong Z(a)\times\mathbb R \text{ when complete},\\
    \text{singular set} \longrightarrow Z_{\mathrm{reg}}(a),\,Z_{\mathrm{sing}}(a).
\end{gathered}
\]
This is a foundation for later developments, not a claim that every arithmetic
expression, singularity, or computational phenomenon is already captured by
these levels.

\subsection{What moved to Paper 0}

The geometric development that earlier versions of this paper contained is now
proved in Paper~0: the arithmetic motions of an arithmetic expression space, the
basic hyperbolic model $\mathfrak E_0$ and the punctured model
$\mathfrak E_1$, the continuous arithmetic motion and its rectification, the
point-path and ripple-pencil readings of a chronological matrix product, the
affine cocycles and the relative affine defect, the accumulative commutative
space with its weighted torsion--Stokes theorem, the reading of torsion as
tearing, and the contact structure with curvature $\mu\lambda$.  Those results
keep their labels, hypotheses, and statuses; this paper cites them and does not
maintain a second version.  The migration is recorded in
\texttt{governance/00b-paper-0-geometric-foundation-amendment.md}.

\subsection{Interface to Paper II}

\emph{Arithmetic Expression Geometry II: Hyperbolic Real Function Theory}
begins where the horizontal connection of Paper~0 stops.  It develops additional
metric and analytic data: a horizontal metric, compatible complex or
almost-complex structures where appropriate, second-order operators,
boundary-value constructions, and arithmetic holomorphicity.  None of these
structures follows uniquely from the contact form alone, and no result of that
downstream function theory is used in the proofs of this paper
\cite{Yuan2026AEGFunctionTheory}.  The $G_4$ sublanguage supplies an additional,
explicitly arithmetic test case: an automorphic or equivariant function used
there must be derived from the group action and must keep its analytic descent
conditions separate from ordinary arithmetic admissibility of a history at a
real input.  The geometry imported by Paper~II --- the regular-AES definition,
the canonical frame, the basic model, and the contact model --- is imported from
Paper~0.

\subsection{Interface to Paper III}

\emph{Arithmetic Expression Geometry III: Singular Zero Geometry and Tubes}
studies the boundary of the regular-zero theorem.  Its objects include singular
and multiple zeros, parameter discriminants, regular and singular zero surfaces,
and tube constructions.  Braid, monodromy, and knot structures there require
separate invariance and normalization arguments.  The Hecke interface proved
here offers a source-driven alternative to imposing a four-valent cellulation by
hand: Paper~III uses a normalized Hauptmodul to construct the corresponding
singular zero dessin, then relates its sign cover to Hecke periodic-orbit knots
\cite{Yuan2026AEGZeroGeometry}.  Those analytic, singular, and three-dimensional
results remain downstream results, not results of this paper.  The complete
regular splitting theorem shows that such a branching network cannot be the zero
set of one connected, complete, boundaryless regular scalar AES; any
singularity, boundary, or incompleteness used in the construction must therefore
be stated explicitly.  In particular, the existence of a tube does not by itself
produce a knot invariant, and no such general invariant is asserted here.

\subsection{Interface to Paper IV}

\emph{Arithmetic Expression Geometry IV: Projective Condensation and
Computational Complexity} returns to the full projective process--result
hierarchy.  Its scope includes bivaluation, quotient and condensation
constructions, representation complexity, and the relation between geometric and
computational resources.  The marked-history formalism and the affine embedding
in projective semantics proved here supply an entry point, and the affine
torsion that Paper~IV compares against is the one proved in Paper~0.  Paper~IV
proves model-relative quotient and residual results together with fixed-model
resource calibrations, while general machine-robust complexity claims remain
open and require explicit models of representation, input size, and computation.

\subsection{Open foundational problems}

Several questions remain at the level of the history layer.
\begin{enumerate}
    \item Characterize the marked-history fibers of the downstream $G_4$
    operator--zero-network realization and determine what arithmetic decoration,
    if any, survives its operator and knot quotients.
    \item Produce fully verified singular multi-zero models and identify their
    stable local normal forms and discriminants; the single verified
    isolated-zero model is a starting point, not a classification.
    \item Determine when the total zero set of a parameter family is proper ---
    hence a locally trivial tube --- rather than merely a submersive family, and
    state the minimal hypotheses under which that upgrade is possible.
    \item Formulate, and then test, a quantitative relation between history
    complexity, representation complexity, and computational complexity, with an
    explicit state space, encoding, and cost model.
    \item Determine which relations among marked histories are compatible with
    ordinary admissibility, as opposed to projective non-degeneracy alone; the
    Hecke order-four example shows that the two conditions are genuinely
    different.
\end{enumerate}

These are research interfaces rather than consequences of the present paper.
The stable conclusion is narrower and more concrete: sequential arithmetic
histories possess projective operator semantics with a distinguished arithmetic
sublanguage; their affine sector is the Borel sector of that semantics; and the
regular zero geometry of an arithmetic expression space is rigid in exactly the
sense stated above, so that singular and branching phenomena require an explicit
change of category.
